\chapter{Segurança da Informação}
\label{cap:seguranca}

\begin{edital}
Políticas e procedimentos de segurança; ISO/IEC 27001:2022 e 27002:2022;
confidencialidade, integridade e disponibilidade; mecanismos de segurança
(controle de acesso, OAuth2, SSO); gestão de riscos (ameaça, vulnerabilidade,
impacto); SDL e OWASP Top 10; análise estática e dinâmica de código (SAST e
DAST). Também: HTTPS, SSL/TLS.
\end{edital}
\peso{ALTO.}

\section{Tríade CIA (a base de tudo)}

\begin{itemize}
\item \textbf{Confidencialidade:} só quem tem direito acessa/\textbf{vê}
(cifra, controle de acesso).
\item \textbf{Integridade:} dado \textbf{íntegro}, não é alterado
indevidamente (hash, assinatura).
\item \textbf{Disponibilidade:} \textbf{acessível} quando necessário
(redundância, backup).
\item Complementos: autenticidade, não repúdio, legalidade.
\end{itemize}

\begin{pegadinha}
Mapeie o verbo do cenário: ``restringir quem vê'' é confidencialidade, não
integridade; ``ficar no ar'' é disponibilidade; ``alteração indevida'' é
integridade. A FGV troca a associação cenário $\to$ princípio.
\end{pegadinha}

\section{Criptografia}

\begin{center}
\begin{tabular}{@{}p{2.6cm}p{4.6cm}p{4.6cm}@{}}
\toprule
& \textbf{Simétrica} & \textbf{Assimétrica (chave pública)} \\
\midrule
Chaves & uma chave secreta compartilhada & par: pública + privada \\
Velocidade & rápida & lenta \\
Uso & cifrar volume de dados & troca de chave, assinatura \\
Exemplos & AES, DES/3DES & RSA, ECC \\
\bottomrule
\end{tabular}
\end{center}

\begin{itemize}
\item \textbf{Hash} (MD5, SHA-256): resumo de tamanho fixo,
\textbf{unidirecional}; garante \textbf{integridade}, não confidencialidade
(não ``descriptografa''). MD5 e SHA-1 são considerados fracos.
\item \textbf{Assinatura digital:} hash do documento cifrado com a
\textbf{chave privada} do emissor $\to$ garante autenticidade, integridade e
não repúdio.
\item \textbf{Certificado digital / ICP-Brasil:} vincula identidade a uma
chave pública, emitido por uma AC (Autoridade Certificadora).
\item \textbf{TDE (Transparent Data Encryption):} cifra os dados \textbf{em
repouso} (arquivos do banco/disco) de forma transparente à aplicação ---
protege contra roubo do arquivo físico, não contra acesso lógico autorizado.
\item \textbf{Esteganografia} $\neq$ criptografia: \textbf{oculta a
existência} da mensagem (esconde dado dentro de imagem/áudio); a
criptografia oculta o \textbf{conteúdo}, não o fato de haver mensagem.
\end{itemize}

\begin{pegadinha}
Hash não é cifra reversível; confidencialidade vem de cifra (não de hash);
assina-se com a chave \textbf{privada}, verifica-se com a \textbf{pública};
esteganografia esconde a existência, cifra esconde o conteúdo. HMAC (chave
simétrica) garante autenticidade e integridade --- não confidencialidade.
\end{pegadinha}

\begin{conceito}[HMAC e o cenário do \emph{webhook}]
\textbf{HMAC} (\emph{Hash-based Message Authentication Code}) é hash
\textbf{combinado a uma chave secreta compartilhada}. Como só quem tem a
chave produz o código, ele entrega \textbf{integridade} (a mensagem não foi
alterada) \emph{e} \textbf{autenticidade} (veio de quem tem a chave). O que
ele \textbf{não} entrega: \textbf{confidencialidade} --- a mensagem viaja
legível --- nem \textbf{não repúdio}, porque a chave é \emph{simétrica} e as
duas pontas conseguem gerar o mesmo código.

\textbf{O cenário que a FGV monta:} um sistema externo envia um
\emph{webhook} (uma chamada HTTP de notificação) e é preciso garantir que a
requisição veio mesmo do parceiro e não foi adulterada no caminho. A resposta
é assinar o corpo com \textbf{HMAC} usando um segredo combinado entre os dois,
e o receptor recalcular e comparar. Hash puro não serve (qualquer um
recalcula); cifrar o corpo resolveria confidencialidade, que não é o problema.
\end{conceito}

\begin{pegadinha}
As duas armadilhas do HMAC: vendê-lo como garantia de
\textbf{confidencialidade} (não é --- não há cifra do conteúdo) e como
garantia de \textbf{não repúdio} (não é --- para isso é preciso
\textbf{assinatura digital} com chave \emph{privada}, que só o emissor
possui). Guarde o corte: chave \emph{simétrica} $\to$ autenticidade +
integridade; chave \emph{privada} $\to$ acrescenta o não repúdio.
\end{pegadinha}

\section{HTTPS, SSL e TLS}

\textbf{HTTPS = HTTP sobre TLS/SSL.} Garante confidencialidade + integridade
+ autenticação do servidor.

\textbf{TLS substitui o SSL:} corrige vulnerabilidades das versões antigas
do SSL. SSL está obsoleto; TLS 1.2/1.3 é o atual.

\begin{pegadinha}
Nunca ``SSL é mais seguro/moderno que TLS'' (é o contrário: TLS sucedeu e
corrigiu o SSL) nem tratar os dois como intercambiáveis.
\end{pegadinha}

\section{Controle de acesso}

\begin{center}
\begin{tabular}{@{}p{3cm}p{8.2cm}@{}}
\toprule
\textbf{Política} & \textbf{Base da decisão} \\
\midrule
DAC (discricionário) & dono do recurso concede acesso \\
MAC (mandatório) & \textbf{rótulos de segurança} comparados com autorizações; regra central, o usuário não decide \\
RBAC (por papéis) & acesso via papéis/funções \\
ABAC (por atributos) & atributos de usuário/recurso/contexto \\
\bottomrule
\end{tabular}
\end{center}

\textbf{Menor privilégio:} só o acesso necessário (ex.: gerente de RH acessa
só o que precisa).

\textbf{OAuth2:} protocolo de \textbf{autorização} delegada (tokens);
\textbf{diferente} de autenticação. \textbf{OpenID Connect (OIDC)} (sobre
OAuth2) é quem faz autenticação --- claims do ID Token: \textbf{iat} (issued
at, emissão), \textbf{exp} (expiração), \textbf{sub} (subject, identificador
do usuário), \textbf{jti} (id do token).

\textbf{SSO (Single Sign-On):} um login para vários sistemas.

\begin{pegadinha}
MAC = rótulos e classificações impostos pelo sistema (o usuário \emph{não}
decide); DAC = o proprietário concede; RBAC = permissões por função/cargo. A
FGV descreve rótulos comparados a autorizações (isso é MAC) e oferece
``discricionário'' como pegadinha. OAuth2 é autorização, \textbf{não}
autenticação. Nos claims do OIDC, a banca pede um e oferece os outros ---
decore os quatro.
\end{pegadinha}

\section{OWASP Top 10 --- 2025 (vigente) e 2021 (leia os dois)}

O edital aponta para o projeto OWASP sem fixar ano, e hoje existem duas
numerações válidas. A \textbf{edição vigente é a 2025} --- oitava da série,
com \emph{release candidate} apresentado em novembro de 2025 e versão final
publicada em janeiro de 2026. A \textbf{2021} é a que a Dataprev 2024 cobrou e
a que ancora as questões deste material. Saber as duas cobre qualquer recorte
que a prova use; e, ao ler o enunciado, \textbf{verifique qual ano ele cita}
antes de contar posição.

\begin{center}
\begin{tabular}{@{}p{1.3cm}p{5.3cm}p{5.3cm}@{}}
\toprule
\# & \textbf{Top 10:2025} (vigente) & \textbf{Top 10:2021} (Dataprev 2024) \\
\midrule
1 & Broken Access Control (absorveu o SSRF) & Broken Access Control \\
2 & Security Misconfiguration (subiu de 5º) & Cryptographic Failures \\
3 & Software Supply Chain Failures (novo; expande ``componentes vulneráveis'') & Injection (SQLi; \textbf{XSS absorvido aqui}) \\
4 & Cryptographic Failures (caiu de 2º) & Insecure Design (novo em 2021) \\
5 & Injection (caiu de 3º) & Security Misconfiguration \\
6 & Insecure Design & Vulnerable and Outdated Components \\
7 & Authentication Failures & Identification and Authentication Failures \\
8 & Software \emph{or} Data Integrity Failures & Software \emph{and} Data Integrity Failures \\
9 & Security Logging and \textbf{Alerting} Failures & Security Logging and \textbf{Monitoring} Failures \\
10 & Mishandling of Exceptional Conditions (novo) & Server-Side Request Forgery (SSRF) \\
\bottomrule
\end{tabular}
\end{center}

O que mudou de 2021 para 2025: \textbf{duas categorias novas} (A03 Software
Supply Chain Failures e A10 Mishandling of Exceptional Conditions) e
\textbf{uma consolidação} (o SSRF, que era A10 em 2021, foi absorvido pelo
Broken Access Control). Três categorias só \textbf{mudaram de nome}: A07
perdeu o ``Identification and''; A08 trocou o ``and'' pelo ``or''; A09 trocou
\textbf{Monitoring} por \textbf{Alerting}.

\begin{pegadinha}
A Dataprev 2024 descreveu ``injeção'' citando SQL $\to$ categoria
\textbf{Injection}; lembre que desde 2021 \textbf{XSS entrou em Injection}
(era categoria própria até 2017). A banca mistura itens que \textbf{não} são
da lista OWASP web --- ``proteção da cadeia de suprimentos'', ``ambiente de
engenharia'', ``treinamento operacional'' são de CI/CD e SSDF, não OWASP
(atenção: em 2025 a cadeia de suprimentos \emph{de software} virou categoria
própria, a A03 --- o que não torna ``proteção da cadeia de suprimentos'' item
de 2021). As renomeações são material de pegadinha pronta: oferecer
``Security Logging and \textbf{Monitoring}'' num item que diz 2025, ou
``Authentication Failures'' num item que diz 2021.
\end{pegadinha}

\section{Desenvolvimento seguro}

\begin{itemize}
\item \textbf{SDL (Security Development Lifecycle):} segurança em todo o
ciclo.
\item \textbf{SAST (estática):} analisa o \textbf{código-fonte} sem
executar (caixa-branca).
\item \textbf{DAST (dinâmica):} testa a \textbf{aplicação em execução}
(caixa-preta).
\item \textbf{IAST:} combina os dois em runtime instrumentado.
\item \textbf{DevSecOps:} segurança automatizada no pipeline CI/CD.
\end{itemize}

\begin{pegadinha}
\textbf{SAST $\times$ DAST}: código parado $\times$ app rodando. Não
inverta.
\end{pegadinha}

\section{X.800 --- a arquitetura de segurança OSI}

Recomendação da ITU-T que dá o \textbf{vocabulário} de segurança usado por
normas e por bancas. Divide o assunto em três: \emph{ataques}, \emph{serviços}
e \emph{mecanismos}. O que a FGV cobra é a divisão dos \textbf{mecanismos}.

\textbf{Os cinco serviços de segurança:} autenticação, controle de acesso,
confidencialidade, integridade e \textbf{não repúdio} (a disponibilidade
aparece como categoria adicional).

\textbf{Mecanismos --- e o corte que a banca pede:}

\begin{center}
\begin{tabular}{@{}p{3.4cm}p{8.6cm}@{}}
\toprule
\textbf{Específicos} \newline (ligados a uma camada e a um serviço) & cifração (\emph{encipherment}), assinatura digital, controle de acesso, integridade de dados, troca de autenticação, \textbf{preenchimento de tráfego} (\emph{traffic padding}), controle de roteamento e \textbf{notarização} \\
\addlinespace
\textbf{Disseminados} \newline (\emph{pervasive}, não específicos de camada) & funcionalidade confiável, \textbf{rótulos de segurança}, detecção de eventos, \textbf{trilha de auditoria} de segurança e recuperação de segurança \\
\bottomrule
\end{tabular}
\end{center}

\begin{pegadinha}
A troca é sempre entre as duas colunas: oferecer \textbf{trilha de auditoria}
ou \textbf{rótulo de segurança} como mecanismo ``específico'' (são
\textbf{disseminados}), ou \textbf{cifração}/\textbf{notarização} como
``disseminado'' (são \textbf{específicos}). O critério é: mecanismo específico
implementa \emph{um serviço} numa \emph{camada}; disseminado é de gestão, vale
para o sistema inteiro. Note que \textbf{preenchimento de tráfego} e
\textbf{controle de roteamento} costumam surpreender --- são específicos, e
servem à confidencialidade do \emph{fluxo}, não do conteúdo.
\end{pegadinha}

\section{Continuidade de negócio: RTO $\times$ RPO}

Os dois objetivos que o plano de continuidade fixa para cada serviço crítico.
Guardar a distinção pela \textbf{unidade do que se perde}:

\begin{center}
\begin{tabular}{@{}p{1.6cm}p{4.4cm}p{6cm}@{}}
\toprule
\textbf{Sigla} & \textbf{Nome} & \textbf{O que limita} \\
\midrule
\textbf{RTO} & Recovery \textbf{Time} Objective & \textbf{TEMPO}: quanto o serviço pode ficar indisponível até ser restabelecido \\
\textbf{RPO} & Recovery \textbf{Point} Objective & \textbf{DADO}: até que ponto no passado se aceita perder informação \\
\bottomrule
\end{tabular}
\end{center}

É o \textbf{RPO} que determina a \textbf{frequência do backup}: aceitar perder
no máximo 15 minutos de dado obriga a proteger os dados a cada 15 minutos. Já
o RTO cobra da infraestrutura de recuperação (redundância, sítio alternativo,
tempo de restauração).

\textbf{As siglas vizinhas que a banca oferece junto:}
\begin{itemize}
\item \textbf{MTBF} (Mean Time Between Failures): tempo \textbf{médio entre}
falhas sucessivas --- métrica de \emph{confiabilidade}, não objetivo de plano.
\item \textbf{MTTR} (Mean Time To Repair): tempo \textbf{médio de reparo}.
\item \textbf{SLA} (Service Level Agreement): o \textbf{acordo} de nível de
serviço --- compromisso contratual, não a métrica em si.
\end{itemize}

\begin{pegadinha}
O cenário dá dois números na ordem ``pode ficar fora X, pode perder Y'' e a
alternativa \textbf{inverte as siglas} (RPO de 4 horas, RTO de 15 minutos).
Ancore em uma letra: \textbf{T} de RTO é \textbf{Tempo}; \textbf{P} de RPO é
\textbf{Ponto} no tempo (\emph{dado}). Os distratores de MTBF/MTTR/SLA
descrevem médias observadas e compromisso contratual --- nenhum é objetivo
definido por diretoria em plano de continuidade. Backup incremental $\times$
diferencial está no Cap.~\ref{cap:orfaos}.
\end{pegadinha}

\section{Detecção e resposta}

\textbf{IDS} (detecção): \textbf{alerta} sobre intrusão, passivo. \textbf{IPS}
(prevenção): \textbf{bloqueia} ativamente (inline). \textbf{SIEM:}
correlaciona logs/eventos para detecção e resposta.

\textbf{Resposta a incidentes (NIST SP 800-61) --- a ordem cai:}
\begin{enumerate}
\item \textbf{Preparação} (antes de tudo: time, ferramentas, políticas).
\item \textbf{Detecção e Análise} (identificar e entender o incidente).
\item \textbf{Contenção, Erradicação e Recuperação} (isolar, remover a
causa, restaurar serviços/backups).
\item \textbf{Atividade pós-incidente / Lições Aprendidas} (documentar,
prevenir recorrência --- realimenta a Preparação).
\end{enumerate}

\begin{pegadinha}
Depois de conter e erradicar vem \textbf{recuperação + lições aprendidas} (a
etapa que ``fecha'' o ciclo); detecção e análise vêm \textbf{antes} da
contenção --- a FGV embaralha essa ordem.
\end{pegadinha}

\begin{jacaiu}
\textbf{Em prova real da FGV:} é o bloco com mais questões reais dos
específicos (38), e a lista abaixo é quase toda verdadeira. Controle de acesso
\textbf{mandatório} (comparação de rótulos); \textbf{OWASP Top 10:2021},
identificar a categoria válida (gabarito: SSRF, o A10 --- as demais
alternativas descreviam \emph{práticas}, não vulnerabilidades); \textbf{X.800}
(mecanismo específico $\times$ disseminado); \textbf{SSL $\times$ TLS} ---
\textbf{Dataprev 2024}. Tríade \textbf{CIA} em associação de colunas; modelo
de \textbf{responsabilidade compartilhada} (PaaS); \textbf{HMAC} em webhook; o
\emph{claim} \texttt{sub} do OIDC para rastreabilidade; aplicabilidade da LGPD
a empresa privada --- \textbf{TJ-RJ}. Anonimização como tratamento que rompe a
associação ao titular; tipo de controle na ISO 27001; o \emph{claim}
\texttt{iat} do OIDC --- \textbf{MPU}. \textbf{RTO $\times$ RPO} saindo de uma
BIA; \textbf{IDS $\times$ IPS} (o que \emph{age}, não só detecta); SSL/TLS;
responsabilidade compartilhada em SaaS; \textbf{RBAC}; segregação de funções;
proprietário do ativo na 27002; tratamento de risco e risco residual; matriz
5$\times$5 com WAF; criptografia simétrica $\times$ assinatura digital sobre
ICP; a família de \emph{malware} inteira (\emph{worm}, \emph{ransomware},
\emph{spyware} + \emph{keylogger}, \emph{phishing}); DDoS e SYN
\emph{flood}; \emph{spoofing} de MAC/IP; NGFW, WAF, proxy e VPN; e as etapas
de \textbf{resposta a incidentes} (contenção, depois recuperação e lições
aprendidas) --- \textbf{ALERO 2026}, que sozinha respondeu por 17 questões de
segurança.

\textbf{No nosso banco} (previsto pelo edital, ainda não visto na amostra de
provas): \textbf{SAST $\times$ DAST} --- nenhuma das 432 questões reais cita
qualquer um dos dois. Continua valendo o estudo: é par de edital e a FGV adora
inverter estático com dinâmico.

Rode \texttt{./quiz.py seguranca}.
\end{jacaiu}

\section{Alta probabilidade / pesquisa extra}

\begin{itemize}
\item \textbf{ISO/IEC 27002:2022:} \textbf{93 controles} em \textbf{4
temas} --- organizacionais (37), pessoas (8), físicos (14), tecnológicos
(34). A 27001:2022 é o \textbf{SGSI} (requisitos, Anexo A); a 27002 é o
\textbf{guia de controles}. Trabalho remoto = pessoas; mídia de
armazenamento = física; criptografia = tecnológica; políticas =
organizacional.
\item \textbf{NIST Cybersecurity Framework 1.1:} funções Identify, Protect,
Detect, Respond, Recover (o 2.0 acrescentou \textbf{Govern}).
\item \textbf{Gestão de risco (ISO 27005/31000):} risco = f(ameaça,
vulnerabilidade, impacto); tratamento: mitigar, transferir, aceitar, evitar.
\item \textbf{STRIDE} (modelagem de ameaças): Spoofing, Tampering,
Repudiation, Information disclosure, Denial of service, Elevation of
privilege.
\end{itemize}

\begin{comosair}
CIA: Confidencialidade = quem VÊ; Integridade = dado ÍNTEGRO; Disponibilidade
= ACESSÍVEL. Controle de acesso, decore a frase-chave: MAC = rótulos impostos
pelo sistema; DAC = o proprietário concede; RBAC = permissões por
função/cargo. OWASP Top 10:2021, âncoras seguras: A01 Broken Access Control,
A03 Injection, A10 SSRF. Se a opção falar em ``cadeia de
suprimentos/pipeline/treinamento'', é distrator de outro framework. Gatilho
geral: ``sempre isento'', ``SSL mais seguro'', claim trocado --- releia com
calma.
\end{comosair}
